\section{Validasi Eksternal}

Seluruh pengujian utama menggunakan dataset hasil perluasan generator ABAC Lab. 
Validasi eksternal dilakukan menggunakan \textit{Smart Building IoT} dari \textcite{diaz-rodriguez_abac_2025} untuk menilai apakah temuan tetap konsisten pada data yang dibangkitkan dengan generator, kebijakan, dan distribusi yang independen dari penelitian ini.
Dataset tersebut memiliki rasio \textit{permit} sekitar $0,75$, sehingga kebijakan trivial yang mengizinkan seluruh permintaan sudah menghasilkan \textit{F1-score} $0,858$ dan digunakan sebagai \textit{baseline} minimum yang relevan. 
Karena dataset asli berukuran sekitar 5,7 juta permintaan, data disubsampel secara terstratifikasi menjadi 8.000 entri dengan dimensi efektif 428. 
Hasil pengujian kualitas H1 disajikan pada Tabel \ref{tab:validasi-iot}.

\begin{longtable}{|>{\raggedright\arraybackslash}p{2cm}|>{\raggedright\arraybackslash}p{1.5cm}|>{\raggedright\arraybackslash}p{1.5cm}|>{\raggedright\arraybackslash}p{1.5cm}|>{\raggedright\arraybackslash}p{1.5cm}|>{\raggedright\arraybackslash}p{3.5cm}|}
\caption{Validasi Eksternal Smart Building IoT --- Kualitas ($D=428$; garis sepele F1=0{,}858; 5 seed)}
\label{tab:validasi-iot} \\
\hline
\textbf{Algoritma} & \textbf{F-score} & \textbf{Presisi} & \textbf{Cakupan} & \textbf{$>$sepele} & \textbf{Keterangan} \\
\hline
\endfirsthead

\multicolumn{6}{c}{{\bfseries \tablename\ \thetable{} -- lanjutan dari halaman sebelumnya}} \\
\hline
\textbf{Algoritma} & \textbf{F-score} & \textbf{Presisi} & \textbf{Cakupan} & \textbf{$>$sepele} & \textbf{Keterangan} \\
\hline
\endhead

\hline
\multicolumn{6}{r}{Lanjut ke halaman berikutnya} \\
\endfoot

\hline
\endlastfoot

% ==================== DATA BARIS ====================
BOA & 0{,}998 & 1{,}000 & 0{,}747 & 5/5 seed & Kebijakan hampir dipulihkan sempurna \\
\hline
ACO & 0{,}994 & 1{,}000 & 0{,}742 & 5/5 seed & Setara BOA \\
\hline
compact-PSO & 0{,}964 & 1{,}000 & 0{,}704 & 4/5 seed & Di atas garis sepele \\
\hline
compact-GA & 0{,}961 & 1{,}000 & 0{,}697 & 5/5 seed & Di atas garis sepele \\
\hline
iBOA & 0{,}943 & 1{,}000 & 0{,}674 & 4/5 seed & Setara CoDA-EDA \\
\hline
CoDA-EDA & 0{,}938 & 1{,}000 & 0{,}665 & 5/5 seed & Melampaui sepele di seluruh seed \\
\hline
UBK & 0{,}834 & 1{,}000 & 0{,}539 & 2/5 seed & Satu-satunya yang tak konsisten \\
\hline

\end{longtable}

Validasi eksternal H1 menunjukkan bahwa CoDA-EDA melampaui \textit{baseline} trivial pada seluruh lima \textit{seed}, dengan \textit{F1-score} rata-rata $0,938$ dibandingkan $0,858$. 
Lima algoritma lain menunjukkan pola serupa, sedangkan UBK hanya melampaui \textit{baseline} pada dua dari lima \textit{seed}. 
Uji Friedman menunjukkan perbedaan signifikan antarkondisi $(p=1{,}6\times10^{-3})$. 
Seluruh algoritma mencapai presisi $1,000$ sehingga perbedaan kualitas terutama ditentukan oleh \textit{coverage}. 
CoDA-EDA berada pada tingkat yang sebanding dengan iBOA dan sedikit di bawah keluarga \textit{compact}, BOA, serta ACO, konsisten dengan pola yang sebelumnya ditemukan pada dataset berdimensi menengah.

H2–H4 juga diuji pada dataset ini. 
Untuk H2, pola jejak memori antaralgoritma konsisten dengan hasil ABAC Lab. 
Untuk H3 dan H4, lima jendela diuji pada kondisi stasioner dan kondisi \textit{drift}. 
Pada kondisi stasioner, lantai derau JSD sebesar $0,0065$ berada di bawah $\tau=0{,}01$, sehingga pemicu pembaruan struktur hanya aktif pada satu dari lima jendela. 
Akibatnya, \textit{warm-start} lebih cepat daripada strategi \textit{always} (6,7 dibanding 7,9 detik) dengan kualitas yang serupa (F1 $0,977$ dibanding $0,973$). 
Hasil ini memberikan dukungan tambahan terhadap H3, yang sebelumnya tidak dapat ditunjukkan pada \textit{e-document} akibat lantai derau yang lebih tinggi.

Pada kondisi \textit{drift} yang dibentuk melalui pengurutan berdasarkan atribut lokasi, pemicu struktur aktif pada seluruh lima jendela sebagaimana diharapkan. 
H4 juga konsisten pada kedua kondisi: \textit{warm-start} sekitar 2$,1$–$2,5$ kali lebih cepat daripada \textit{re-mining} penuh (6,7 dibanding 16,9 detik pada salah satu konfigurasi) serta menghasilkan F1 lebih tinggi, yaitu $0,977$ dibanding $0,914$.

Secara keseluruhan, validasi eksternal menunjukkan bahwa pola H1–H4 tetap muncul pada dataset sintetis yang dibangkitkan oleh \textit{generator} independen. 
Hasil ini mendukung ketahanan lintas-\textit{generator}, tetapi tidak dapat ditafsirkan sebagai bukti generalisasi terhadap data operasional \textit{edge-IoT}. 
Khusus H3, keberhasilan mekanisme pada \textit{Smart Building IoT} dengan lantai derau rendah menunjukkan bahwa keterbatasan yang ditemukan pada \textit{e-document} lebih berkaitan dengan karakteristik estimasi distribusi dataset daripada kegagalan mekanisme pembaruan struktur itu sendiri.